\documentclass[11pt,a4paper]{article}
\usepackage[T1]{fontenc}
\usepackage[utf8]{inputenc}
\usepackage{lmodern}
\usepackage[margin=25mm,headheight=14pt]{geometry}
\usepackage{microtype}
\usepackage{enumitem}
\usepackage{booktabs,tabularx}
\usepackage{xcolor}
\usepackage{fancyhdr}
\usepackage[hidelinks]{hyperref}
\definecolor{dsrlgreen}{HTML}{176B52}
\definecolor{dsrlblue}{HTML}{245A73}
\usepackage[colorlinks=true,linkcolor=dsrlgreen,urlcolor=dsrlblue]{hyperref}
\setlist{itemsep=3pt,topsep=5pt}
\setlength{\parskip}{5pt}
\setlength{\parindent}{0pt}
\pagestyle{fancy}
\fancyhf{}
\fancyhead[L]{\small\textsc{Digital Sustainability Readiness Level}}
\fancyhead[R]{\small Initial Manifesto · v0.1}
\fancyfoot[C]{\thepage}
\renewcommand{\headrulewidth}{0.4pt}

\title{\vspace{-1.4cm}\Huge\bfseries\color{dsrlgreen}The DSRL Manifesto\\[0.4em]
\Large\normalfont\color{black}A first statement on digital sustainability readiness}
\author{Rafael Oliveira Rodrigues\\\small \url{https://digisustainability.io/} · \href{mailto:ror@digisustainability.io}{ror@digisustainability.io}}
\date{\small Initial discussion draft · October 2026}

\begin{document}
\maketitle
\begin{center}
\color{dsrlblue}\large\itshape Digital progress should advance human and planetary well-being, not make either an afterthought.
\end{center}
\vspace{0.4cm}

\begin{abstract}
The Digital Sustainability Readiness Level (DSRL) is a proposed framework for helping organizations consider digital transformation and sustainability together. It asks two connected questions: how can digital capabilities contribute to positive environmental and social outcomes, and how can the impacts of digital systems themselves be understood and reduced? This initial manifesto sets out the framework's purpose, principles, seven assessment domains, and nine-level progression. It is a starting point for discussion and validation, not a completed standard, certification scheme, or claim of demonstrated impact. Its next step is collaborative refinement through evidence, piloting, and open critique.
\end{abstract}

\section*{Why a readiness framework?}
Organizations increasingly rely on digital systems while facing pressure to improve environmental performance, social outcomes, and accountability. These agendas are often managed separately. When they are disconnected, a digital initiative may overlook its own energy, material, or social footprint; a sustainability initiative may lack the data, infrastructure, or organizational capability needed to deliver and evaluate change.

DSRL proposes a shared language for examining this relationship. Readiness is not simply the presence of technology or the publication of a sustainability commitment. It concerns the capabilities, evidence, governance, and learning practices that allow an organization to use digital means responsibly and to improve over time.

\section*{Our starting propositions}
\begin{enumerate}
  \item \textbf{Digital and sustainability decisions belong in the same conversation.} Technology choices have environmental and social consequences; sustainability goals can also shape better digital choices.
  \item \textbf{Measure both contribution and footprint.} An assessment should consider where digital solutions may enable sustainability outcomes and where digital operations create impacts or risks.
  \item \textbf{Readiness is multidimensional.} Strong performance in one area does not erase gaps in another. Infrastructure, governance, data, business models, people, and innovation all matter.
  \item \textbf{Progress is contextual and iterative.} A level is a guide to capabilities and evidence, not a judgement of an organization's worth. Organizations should be able to identify a next step and revisit it as conditions change.
  \item \textbf{Claims need evidence.} Scores and public statements should be traceable to appropriate evidence, with assumptions, boundaries, uncertainty, and limitations made visible.
  \item \textbf{The framework must remain open to challenge.} DSRL should be tested with organizations and affected stakeholders, revised in light of evidence, and clear about what remains unsettled.
\end{enumerate}

\section*{Seven connected domains}
The proposed framework examines organizational or project readiness across seven domains. They are connected lenses rather than independent silos.
\begin{enumerate}
  \item \textbf{Environmental Impact \& Resource Efficiency} — understanding and improving environmental impacts associated with digital activity, including energy, emissions, materials, and resource use.
  \item \textbf{Digital Infrastructure Sustainability} — considering the sustainability of digital infrastructure and its lifecycle, including systems, services, procurement, and operational efficiency.
  \item \textbf{Data \& Analytics for Sustainability} — using relevant, governed data and analytics to inform sustainability decisions, monitor outcomes, and learn from results.
  \item \textbf{Sustainable Business Model Innovation} — aligning digital-enabled value creation with durable environmental and social goals, while examining incentives and possible unintended effects.
  \item \textbf{Stakeholder Engagement \& Social Impact} — involving affected stakeholders, considering accessibility and inclusion, and evaluating social outcomes of digital initiatives.
  \item \textbf{Governance \& Compliance} — establishing accountability, decision rights, risk management, transparency, and processes for meeting applicable obligations.
  \item \textbf{Innovation \& Circular Economy Integration} — applying digital capabilities to support circular practices and responsible innovation, while considering lifecycle impacts and collaboration across value chains.
\end{enumerate}

\section*{A nine-level progression}
Each domain is considered on a proposed scale from 1 to 9. The scale describes a broad progression from limited awareness and ad-hoc activity toward integrated, adaptive, and transformative practice. It should not be read as a guarantee of impact or a universal ranking.

\begin{center}
\small
\begin{tabularx}{\textwidth}{>{\bfseries}c l X}
\toprule
Level & Working label & Indicative progression \\
\midrule
1 & Unaware / traditional & No established recognition or practice in the domain. \\
2 & Basic awareness & Initial understanding; actions remain informal or isolated. \\
3 & Initial measurement & Early data, review, or assessment begins to establish a baseline. \\
4 & Strategic planning & Goals, responsibilities, and plans are defined. \\
5 & Pilot implementation & Approaches are tested in bounded settings and reviewed. \\
6 & Scaling and integration & Practices with evidence of value are expanded and integrated into operations. \\
7 & Systematic practice & Responsibilities, processes, and monitoring are embedded across the organization. \\
8 & Continuous improvement & Evidence and feedback support ongoing adaptation and optimization. \\
9 & Transformative leadership & The organization demonstrates learning and leadership beyond its own boundaries. \\
\bottomrule
\end{tabularx}
\end{center}

These labels are an initial synthesis. Domain-specific indicators and the evidence required at each level must be developed and tested before the scale can support consistent comparison. An organization may be at different levels across domains; a single average should not hide that profile.

\section*{How DSRL should be used}
We envisage DSRL as a structured self-assessment and learning aid. A responsible assessment would define its scope, gather relevant evidence, involve people with different responsibilities and perspectives, describe uncertainty, and turn findings into prioritized actions. A domain profile may help reveal strengths and gaps; a roadmap may help teams decide what to test, resource, and review next.

Any future scoring method should explain how ratings are assigned, how evidence is evaluated, whether domains are weighted, and how uncertainty is represented. Benchmarks should only be introduced when data are comparable and their limitations are clear. Certification or external assurance would require separate governance, criteria, and validation; this manifesto does not establish or promise such a scheme.

\section*{Commitments for the next iteration}
To make DSRL useful and credible, its development should:
\begin{itemize}
  \item publish clear definitions, assessment boundaries, and version history;
  \item develop observable, domain-specific indicators for all nine levels;
  \item pilot the assessment in varied organizational contexts, including smaller organizations;
  \item test whether different assessors can apply the criteria consistently;
  \item include affected stakeholders and examine accessibility, equity, and unintended consequences;
  \item distinguish measured outcomes from aspirations, estimates, and illustrative examples; and
  \item revise the framework openly in response to evidence and critique.
\end{itemize}

\section*{A call to collaborate}
We invite researchers, practitioners, public institutions, civil society, and organizations of different sizes to help shape DSRL. Contribute evidence, question assumptions, test the domains, improve the indicators, and share lessons from implementation. The aim is not to add another label to existing practice, but to support better decisions about how digital change is designed, governed, and evaluated.

\begin{center}
\vspace{0.5em}
\Large\bfseries\color{dsrlgreen}Make digital progress accountable to the future it helps create.
\end{center}

\section*{About this draft}
After 1 year of review, this is the first manifesto --- produced with love at Impact Hub Zurich.
This initial manifesto was prepared as a synthesis of the DSRL materials supplied by the author, including the DSRL comprehensive assessment framework (v0.1) and two presentation drafts. It preserves the seven domains and nine-level approach described in those materials while framing DSRL as a proposal for further validation. It does not independently verify quantitative claims, certify organizations, or represent a formally adopted standard. Feedback and subsequent revisions should be recorded in the repository history.

\vfill
\begin{center}\small\color{gray}DSRL Manifesto · Version 0.1 · Discussion draft\end{center}
\end{document}
